\documentclass[10pt,a4paper]{amsart}
\usepackage[margin=19mm]{geometry}
\usepackage{amsmath,amssymb,mathtools}
\usepackage[scr=rsfs,cal=euler]{mathalfa}
\usepackage{xfrac}
\usepackage[unicode,hidelinks,bookmarks=false]{hyperref}
\newcommand{\ie}{i.e.\ }
\newcommand{\cf}{cf.\ }
\newcommand{\resp}{respectively\ }
\newcommand{\Subsection}{\subsection}
\providecommand{\nobreakdash}{\nobreak\hbox{-}}
\setlength{\emergencystretch}{2em}
\multlinegap0pt
\usepackage{mathtools}
\usepackage{amssymb}
\usepackage[shortlabels]{enumitem}
\usepackage{caption}
\usepackage{float}
\usepackage{xcolor}
\usepackage{tcolorbox}
\newtheorem{lemma}{Lemma}
\newtheorem{remark}{Remark}
\title{Asymptotic behavior for a class of damped second-order gradient systems via Lyapunov method}
\author{Renan J. S. Isneri, Eric B. Santiago, Severino H. da Silva}
\date{Source: arXiv:2512.20751v1 (CC BY 4.0)}
\begin{document}
\maketitle

\noindent\emph{Source and licence.} Asymptotic behavior for a class of damped second-order gradient systems via Lyapunov method. Version arXiv:2512.20751v1 (CC BY 4.0), distributed by arXiv under the Creative Commons Attribution 4.0 International licence, http://creativecommons.org/licenses/by/4.0/, as recorded in the arXiv licence metadata for this exact version and shown on the abstract page https://arxiv.org/abs/2512.20751v1. Full licence notice: \texttt{licenses/arxiv-2512.20751-CC-BY-4.0.txt}.\par\smallskip

\noindent\textit{Editorial scope.} Counted unit: the article's lemma L5 (source0.tex lines 534--544). The article's complete original proof of it is delivered with it (source0.tex lines 545--641, one \texttt{proof} environment of 97 lines). No mathematics is written, repaired or re-derived: the statement and the proof are the article's own lines, in the article's own notation, and no retained line is substituted or re-flowed.  Retained uncounted as the source-local context the proof needs: lines 206--208; lines 348--350; lines 378--383; lines 385--387; lines 390--399; lines 440--445; lines 495--501; lines 505--507.  Nothing was omitted from any retained span, so the package carries no omission record.  No retained line contains a citation, so the wrapper carries no bibliography and no bibliography entry.  No retained line contains a figure environment, an image-inclusion command or drawing code, so no image is delivered and the package declares no asset dependency.  Editorial formatting only, in addition: an AMS \texttt{amsart} preamble that restates the article's own declarations and macros used by the retained lines, namely the environments \texttt{lemma}, \texttt{remark} on separate counters, together with the standard packages those lines need.  The retained environments print in this document as Lemma 1, Remark 1 in order of appearance, whereas the article prints the same items with the numbering of its own full document.  The complete native source of the article as distributed in the arXiv e-print for this exact version is bundled unchanged as \texttt{sources/191594/source0.tex}.
\begin{equation}\label{Equation}
	\ddot{u}(t) + a \dot{u}(t) + \nabla W(u(t)) = 0,
\end{equation}

\begin{lemma}\label{L1}
	Every classical solution $u$ of \eqref{Equation} is globally defined.
\end{lemma}

\begin{equation}\label{S}
	\left\{\begin{array}{ll} 
		\dot{x} = y, & \\ 
		\dot{y} = -a y - \nabla W(x). & 
	\end{array}\right.
\end{equation}

\begin{equation}\label{Va}
	V_a(x,y) := \frac{1}{2} \|y\|^2 + 2W(x) + \frac{1}{2} \|y + a x\|^2,
\end{equation}

\begin{lemma}\label{L2}
	The functional $V_a$ is a strict Lyapunov function for system \eqref{S} at the equilibrium $(0,0)$. More precisely, $V_a(0,0)=0$, and for all $(x,y)\in B_\eta(0)\times \mathbb{R}^N$ with $(x,y)\neq(0,0)$ one has
	$$
	0 < V_a(x,y)\quad\text{and}\quad\langle \nabla V_a(x,y), F_a(x,y) \rangle_* < 0,
	$$
	where $F_a(x,y)$ is the vector field of system \eqref{S}, which is given by
	\begin{equation}\label{VectorField}
		F_a(x,y) := (y, -ay - \nabla W(x)).
	\end{equation}
\end{lemma}

\begin{lemma}\label{L3}
Let $m$, $a_0$ and $\kappa$ be positive real numbers. Then, the point $ (0,0) \in \mathbb{R}^{2N} $ is an interior point of the set
$$
\mathcal{W} := \left\{ (x,y) \in B_\kappa(0,0) : V_a(x,y) < m \text{ for all } a \in (0, a_0] \right\}.
$$
\end{lemma}

\begin{remark}\label{OBS2} \rm
If there exists $t_0\in\mathbb{R}$ such that $\varphi_a(t_0;z_0)=(0,0)$, then $\varphi_a(t;z_0)=(0,0)$ for all $t\in \mathbb{R}$. Indeed, by \hyperlink{W1}{$(W_1)$}, $F$ is locally Lipschitz, and so, we have that by the Picard-Lindelöf Theorem, the Cauchy problem
$$
\phi'(t)=F(\phi(t)),\qquad \phi(t_0)=(0,0),
$$
admits a unique solution. Since $F(0,0)=(0,0)$, the constant function $\phi\equiv(0,0)$ is a solution, by uniqueness, $\varphi_a(t;z_0)=(0,0)$ for any $t\in \mathbb{R}$.
\end{remark}

\begin{lemma}\label{L4}
	If $I\subset\mathbb{R}$ is an interval such that $\varphi_a(t;z_0)\in B_\eta(0)\times \mathbb{R}^N$ for all $t\in I$, then the function $t\mapsto V_a(\varphi_a(t;z_0))$ is non-increasing on $I$.
\end{lemma}

\begin{lemma}\label{L5}
Given any $\epsilon > 0$ and $a_0>0$, there exists $\sigma > 0$ such that, for every $a\in (0,a_0]$ and for every solution $\varphi_a(t;z_0)$ of system \eqref{S}, corresponding to the damping parameter $a$, satisfying
$$
\|z_0\|_* < \sigma,
$$
we have
$$
\|\varphi_a(t;z_0)\|_* < \epsilon, \quad \text{for all } t \geq 0,\quad \text{and}\quad \lim_{t \to +\infty} \varphi_a(t;z_0) = (0,0).
$$
In other words, the asymptotic stability of the equilibrium (0,0) is uniform with respect to the damping parameter.
\end{lemma}

\begin{proof}
For each $a>0$ let $V_a:B_\eta(0)\times \mathbb{R}^N\to \mathbb{R}$ the Lyapunov function for $(0,0) \in \mathbb{R}^{2N}$ defined as in \eqref{Va}. Now, let $\epsilon > 0$ such that
$$
\overline{B_\epsilon(0,0)} = \left\{ (x,y) \in \mathbb{R}^{2N} : \|(x,y)\|_* \leq \epsilon \right\} \subset B_\eta(0)\times \mathbb{R}^N.
$$
We consider the following set 
$$
M = \left\{ V_a(x,y) : \|(x,y)\|_* = \epsilon \text{ and } a>0 \right\}.
$$
We observe that $M$ is bounded below in $\mathbb{R}$. Indeed, just note that $V_a(x,y)\ge 0$ for all $a>0$ and $(x,y)\in B_\eta(0)\times \mathbb{R}^N$. Thereby, 
$$
m := \inf M \geq 0.
$$
We claim that $m > 0$. Indeed, since 
$$
V_a(x,y)\ge W(x)\,\text{for all } a>0\text{ and }(x,y)\in M,
$$
and $W$ is continuous, as $W$ is restricted to a compact subset of $B_\eta(0)$, we conclude by \hyperlink{W1}{$(W_1)$} that $m > 0$. Now, we define the set
$$
\mathcal{W}= \left\{ (x,y) \in B_\epsilon(0,0) : V_a(x,y) < m \text{ for all }a\in (0,a_0]\right\}.
$$
From Lemma \ref{L3}, $(0,0)$ is an interior point of $\mathcal{W}$. So, there exist $\sigma  \in (0, \epsilon)$ such that $B_\sigma(0,0) \subset \mathcal{W}$.
	
\noindent\textbf{\hypertarget{C1}{Claim 1:}} If $\varphi_a(t;z_0)$ is a solution of system \eqref{S}, corresponding to the damping parameter $a\in (0,a_0]$, with initial condition $ \varphi_a(0;z_0)=z_0\in B_\sigma (0,0)$, then
$$
\|\varphi_a(t;z_0)\|_* < \epsilon \quad \text{for all } t \geq 0.
$$

To see this, note first that by Lemma \ref{L1} we can conclude that $\varphi_a$ is globally defined. Now, we assume by contradiction that Claim \hyperlink{C1}{1} is false. Then, for some $t_{0} > 0$, we must have $\varphi_a(t_{0};z_0) \notin B_\epsilon(0,0)$. Since $z_0 \in B_\epsilon(0,0)$, by continuity, there exists $t_* \in (0, t_{0})$ such that
\begin{equation}\label{3eq}
	\varphi_a(t;z_0) \in B_\epsilon(0,0)\,\text{ for all } \, t \in (0, t_*)\text{ and }\|\varphi_a(t_*;z_0)\|_* = \epsilon.
\end{equation}
Thus, $\varphi_a(t_*;z_0) \in \partial B_\epsilon(0,0)$ and hence not in $\mathcal{W}$. Consequently,
\begin{equation}\label{ineq1}
	V_a(\varphi_a(t_*;z_0)) \geq m > V_a(\varphi_a(0;z_0))=V(z_0).
\end{equation}
On the other hand, from \eqref{3eq}, we can apply Lemma \ref{L4} to ensure that
$$
V_a(\varphi_a(t_*;z_0)) \leq V_a(z_0),
$$
contradicting inequality \eqref{ineq1}. This contradiction implies that $\varphi_a(t;z_0) \in B_\epsilon(0,0)$ for all $t \geq 0$, proving Claim \hyperlink{C1}{1}.
	
\noindent\textbf{\hypertarget{C2}{Claim 2:}} If $z_0 \in B_ \sigma(0,0)$, then
$$
\lim_{t \to +\infty} \varphi_a(t;z_0) = (0,0).
$$
	
First, we observe that by Remark \ref{OBS2} we can assume without loss of generality that $\varphi_a(t;z_0) \neq (0,0)$ for any $t\ge 0$. On the other hand, from Claim \hyperlink{C1}{1}, the solution $\varphi_a(t;z_0)$ remains in $B_\epsilon(0,0)$ for all $t \geq 0$. Let us consider the function again
$$
h(t) := V_a(\varphi_a(t;z_0)), \quad t \geq 0.
$$
We have already discussed in Lemma \ref{L4} that $h \in C^1([0, \infty))$ and, because of Lemma \ref{L2}, 
$$
\dot{h}(t) = \langle \nabla V_a(\varphi_a(t;z_0)), F_a(\varphi_a(t;z_0)) \rangle_* < 0, \, \text{ for all }t > 0,
$$
from which it follows that $h$ monotonically decreasing. Moreover, $h$ is bounded below. Therefore, there exists $d \geq 0$ satisfying
$$
V_a(\varphi_a(t;z_0))\ge d \, \text{ for any }\, t\ge 0\, \text{ and }\, \lim_{t \to +\infty} V(\varphi_a(t;z_0)) = d.
$$
We claim that $d = 0$. Suppose by contradiction that $d > 0$. By continuity of $V$, and the fact that $V(0,0) = 0$, there exists $r \in (0, \epsilon)$ such that
$$
V_a(z) < d \quad \text{for all } z \in B_r(0,0).
$$
Thus, $\varphi_a(t;z_0) \notin B_r(0,0)$ for all $t \geq 0$. Since the function $f(z) = \langle \nabla V_a(z), F_a(z) \rangle_*$ is continuous and $V_a$ is strict Lyapunov functional at the origin, and the set $\overline{B_\epsilon(0,0)} \setminus B_r(0,0)$ is compact, there exists $b > 0$ such that
$$
\langle \nabla V_a(z), F_a(z) \rangle_* \leq -b \quad \text{for all } z \in \overline{B_\epsilon(0,0)} \setminus B_r(0,0).
$$
As $\varphi_a(t;z_0) \in \overline{B_\epsilon(0,0)} \setminus B_r(0,0)$, we get
$$
\dot{h}(t) \leq -b \quad \text{for all } t \geq 0.
$$
Integrating, we obtain
$$
V_a(\varphi_a(t;z_0)) - V_a(z_0) \leq -bt,\, \text{ for all }t>0,
$$
which results in
$$
V_a(\varphi_a(t;z_0)) \to -\infty \quad \text{as } t \to +\infty,
$$
which is a contradiction. Therefore, $d = 0$. To finish, take any sequence $t_n \to +\infty$. Since $\varphi_a(t_n;z_0) \in B_\epsilon(0,0)$ for any $n\in\mathbb{N}$, by Bolzano-Weierstrass, there exists a subsequence $t_{n_k}$ such that
$$
\varphi_a(t_{n_k};z_0) \to w \in \overline{B_\epsilon(0,0)}.
$$
By continuity of $V_a$,
$$
\lim_{k \to +\infty} V_a(\varphi_a(t_{n_k};z_0)) = V_a(w) = 0,
$$
which implies in $w = (0,0)$, because $V_a(z) > 0$ for all $z \neq (0,0)$. Hence,
$$
\lim_{k \to+ \infty} \varphi_a(t_{n_k};z_0) = (0,0).
$$
As every sequence $t_n \to +\infty$ has a subsequence converging to $(0,0)$, it follows that
$$
\lim_{t \to +\infty} \varphi_a(t;z_0) = (0,0),
$$
completing the proof.
\end{proof}

\end{document}
